%% kapa: exam=1 round=30 year=2019 period=2 form=A subject=appraisal_law subject_ko=감정평가관계법규 seq=2 number=2
%% source: Q-Net 감정평가사 기출문제 (cst00309) artlSeq=5205205 file=extracted/1008/5205205/2019년 제30회 감정평가사 1차시험 문제지/2019년 제30회 감정평가사 1차시험 문제지/2019년도 제30회 감정평가사 1차 2교시 A형.hwp
%% conversion: kordoc 4.24.51 -> tools/gen1.py (auto); review=proofread-sonnet; flags=none
\begin{question}{2}
국토의 계획 및 이용에 관한 법령상 광역도시계획에 관한 설명으로 옳지 않은 것은?
\begin{choices}
  \item 광역계획권이 둘 이상의 인접한 시ㆍ도의 관할구역에 걸쳐 있는 경우 국토교통부장관이 광역계획권을 지정한다.
  \item 광역도시계획에는 광역계획권의 지정목적을 이루는 데 필요한 경관계획에 관한 사항이 포함되어야 한다.
  \item 국토교통부장관은 광역도시계획을 수립하거나 변경하려면 미리 공청회를 열어야 한다.
  \item 광역도시계획을 공동으로 수립하는 시ㆍ도지사는 그 내용에 관하여 서로 협의가 되지 아니하면 공동으로 조정을 신청하여야 한다.
  \item 광역도시계획협의회에서 광역도시계획의 수립에 관하여 조정을 한 경우에는 그 조정 내용을 광역도시계획에 반영하여야 한다.
\end{choices}
\end{question}
\answer{4}
